% ============================================================
%  吴宛幸  个人简历 (中文版)
% ============================================================
\documentclass[10pt, a4paper]{article}

% ── Encoding & CJK Support ──────────────────────────────────
\usepackage[UTF8, heading=false, fontset=none]{ctex}
\setCJKmainfont{Heiti SC}[AutoFakeBold=true, AutoFakeSlant=true]
\setCJKsansfont{Heiti SC}[AutoFakeBold=true]
\setCJKmonofont{Heiti SC}[AutoFakeBold=true]
\usepackage[T1]{fontenc}
\usepackage{palatino}
\usepackage{mathpazo}
\usepackage{microtype}

% ── Page Layout ──────────────────────────────────────────────
\usepackage[top=1.0cm, bottom=1.0cm, left=1.35cm, right=1.35cm]{geometry}
\setlength{\parindent}{0pt}
\setlength{\parskip}{0pt}
\setlength{\emergencystretch}{3em}
\linespread{1.00}\selectfont

% ── Packages ─────────────────────────────────────────────────
\usepackage{xcolor}
\usepackage{hyperref}
\usepackage{enumitem}
\usepackage{titlesec}
\usepackage{tabularx}
\usepackage{array}
\usepackage{fancyhdr}
\usepackage{tikz}
\usepackage{etoolbox}
\usepackage{graphicx}

% ── Colors ───────────────────────────────────────────────────
\definecolor{headcol}{RGB}{30, 30, 30}
\definecolor{rulecol}{RGB}{180, 50, 30}
\definecolor{secline}{RGB}{200, 200, 200}
\definecolor{linkcol}{RGB}{40, 80, 160}
\definecolor{dategray}{RGB}{100, 100, 100}

% Tag colors
\definecolor{tagroute}{RGB}{130, 50, 160}
\definecolor{taggenomics}{RGB}{40, 140, 90}
\definecolor{tagrl}{RGB}{180, 100, 20}
\definecolor{tagautoresearch}{RGB}{30, 100, 160}

\hypersetup{
  colorlinks = true,
  urlcolor   = linkcol,
  linkcolor  = linkcol,
  pdfauthor  = {吴宛幸},
  pdftitle   = {吴宛幸，个人简历},
}

% ── Section Headers ──────────────────────────────────────────
\titleformat{\section}
  {\large\bfseries\color{headcol}}
  {}{0em}{}
  [\vspace{-5pt}{\color{secline}\rule{\linewidth}{0.4pt}}\vspace{0pt}]
\titlespacing*{\section}{0pt}{6pt}{2pt}

% ── Badge Tags ───────────────────────────────────────────────
\newcommand{\badge}[2]{%
  \tikz[baseline=(x.base)]{
    \node[inner xsep=3.5pt, inner ysep=1.6pt, rounded corners=2.5pt,
          draw=#1, line width=0.45pt,
          text=#1, font=\fontsize{6}{7}\selectfont\sffamily]
      (x) {#2};
  }%
}

\newcommand{\troute}  {\badge{tagroute}{routing}}
\newcommand{\tgenomics}{\badge{taggenomics}{genomics}}
\newcommand{\trl}     {\badge{tagrl}{agentic}}
\newcommand{\topd}    {\badge{tagrl}{opd}}
\newcommand{\tautoresearch} {\badge{tagautoresearch}{autoresearch}}

% ── Bullet Lists ─────────────────────────────────────────────
\setlist[itemize,1]{
  leftmargin = 1.3em,
  label      = {\small\textbullet},
  itemsep    = 0pt,
  topsep     = 1pt,
  parsep     = 0pt,
}

% ── Helpers ──────────────────────────────────────────────────
\newcommand{\unilogo}[1]{%
  \raisebox{-0.5pt}{\includegraphics[height=8pt]{assets/#1}}%
}
\newcommand{\me}[1]{\textbf{#1}}
\newcommand{\eqstar}{$^{*}$}
\newcommand{\corr}{$^{\dagger}$}

\def\pubskip{5pt}

% Single publication block
\newcommand{\pub}[4]{%
  \vspace{\pubskip}%
  \noindent\textbf{#1}\ifstrempty{#2}{}{\enskip #2}%
  \begin{itemize}[topsep=0pt]
    \item #3
    \item \textit{#4}
  \end{itemize}%
}

% Publication block with extra "highlight" bullet
\newcommand{\paperentry}[4]{%
  \par\vspace{5pt}%
  {\raggedright\textbf{#1}\ifstrempty{#2}{}{\enskip #2}\par}%
  {\small #3\par}%
  {\small\textit{#4}\par}%
}
\newcommand{\researchitem}[3]{%
  \par\vspace{7pt}%
  {\textbf{#1}\enskip #2\par}%
  \vspace{1pt}#3\par
}

% Education / Experience entry
\newcommand{\entry}[2]{%
  \noindent
  \begin{tabularx}{\linewidth}{@{}X r@{}}
    \textbf{#1} & {\small\color{dategray}#2}
  \end{tabularx}%
  \par\vspace{0pt}%
}

% ── Header / Footer ──────────────────────────────────────────
\pagestyle{empty}

% ============================================================
\begin{document}
\thispagestyle{empty}
\fontsize{10.4}{12.2}\selectfont

% ── NAME BLOCK ───────────────────────────────────────────────
\noindent
\begin{tabularx}{\linewidth}{@{}X r@{}}
  \raisebox{-4pt}{%
    {\fontsize{28}{32}\selectfont\bfseries 吴宛幸}%
  }
  &
  \begin{tabular}[b]{@{}r@{}}
    \small 150 6798 2279 \quad \href{mailto:wuwanxing3574@gmail.com}{wuwanxing3574@gmail.com}\\[2pt]
    \small 微信：wwx15067982279\quad\href{https://github.com/wwx777}{GitHub}\quad\href{https://wwx777.github.io}{Homepage}
  \end{tabular}
\end{tabularx}

\vspace{4pt}
{\color{rulecol}\rule{\linewidth}{1.2pt}}
\vspace{1pt}

\section{教育经历}
\entry{\unilogo{ipp.png}\enskip 巴黎理工学院（Institut Polytechnique de Paris）}{法国巴黎}
\textit{计算机科学（DS4Health 方向）硕士}\hfill\textit{2025.09 -- 2028.05}\\
GPA：\textbf{16.6/20}（A+）；排名：\textbf{前 5\%}。\\
{\small \href{https://www.topuniversities.com/universities/institut-polytechnique-de-paris}{QS Ranking: 43}；\href{https://www.timeshighereducation.com/world-university-rankings/institut-polytechnique-de-paris}{THE CS Ranking: 43}。}

\vspace{4pt}
\entry{\unilogo{sustech.png}\enskip 南方科技大学}{深圳}
\textit{数据科学与大数据技术}\hfill\textit{2021.09 -- 2025.07}\\
本科导师：\href{https://ghchen.me/}{陈冠华}。

\section{科研经历}
\entry{\unilogo{sustech.png}\enskip SUSTech-NLP 实验室}{深圳}
\textit{科研助理}\hfill\textit{2025.01 -- 至今}\\
合作导师：\href{https://ghchen.me/}{陈冠华}。
\researchitem{Agent 探索与经验学习}{\trl}{Agent 要适应陌生环境，需要先探索环境、收集信息，再据此完成任务。本研究关注\textbf{如何高效获取有用的探索经验}，通过两组控制实验分析环境线索、候选空间和检查数量对探索效果的影响。提出 \textbf{Scope-then-Cover}：强 actor 根据任务选定探索范围，弱 scout 并行检查并返回观察。将这一方法用于家务、电商和代码 Agent 任务，帮助 Agent 获取决策所需的信息；再利用探索轨迹开展在线学习、积累经验库，改善后续任务表现。}
\researchitem{多视角 on-policy 自蒸馏}{\badge{tagrl}{post-training}}{研究如何将同一模型在不同条件下展现的\textbf{互补推理能力整合进单一策略}。以跨语言推理为切入点，提出 \textbf{CMOPD 跨语言多视角 on-policy 自蒸馏}：融合冻结模型在不同语言条件下的逐 token 监督，指导学生自身生成的推理轨迹。训练后，模型直接使用目标语言推理，无需多语言生成与答案聚合。比较三个规模的模型在四种低资源语言下的推理表现，分析不同语言视角提供了哪些互补信息，以及如何融合监督才能帮助模型学习。}
\researchitem{端云协同模型路由评测}{\troute}{端云协同让本地小模型处理日常请求，并在需要时借助云端大模型。但云端预算有限，哪些问题应该优先获得帮助？围绕\textbf{如何评价路由器对云端资源的分配}，构建 \textbf{RouterXBench}，同时考察它能否识别小模型难以处理的问题，以及这些选择带来的回答质量和云端成本。比较不同路由方法，讨论\textbf{优先解决困难问题与帮助更多问题之间的取舍}，以及不同预算和质量要求如何影响路由方法的选择。围绕指令遵循、知识与数学组织六项测试，并提供轻量参考路由器 \textbf{MeanProbe}。}

\section{学术论文}
{\small $^*$共同第一作者；$^\dagger$通讯作者。}
\paperentry{Can Weak Models Explore? Scope Then Cover for LLM Agent Adaptation}{\trl}{\me{Wanxing Wu}, Zhiwen Ruan, He Zhu\corr, Guanhua Chen\corr}{ICLR 2027 在投}
\paperentry{Cross-lingual Multi-view On-Policy Self-Distillation: Distilling Multilingual Views into a Single Target-Language Policy}{\badge{tagrl}{post-training}}{\me{Wanxing Wu}, He Zhu\corr, Guanhua Chen\corr}{ICLR 2027 在投}
\paperentry{Towards Fair Evaluation of Routers in Collaborative LLM Systems}{\troute}{\me{Wanxing Wu}\eqstar, He Zhu\eqstar, Yixia Li\eqstar, \ldots, Guanhua Chen\corr}{NAACL 在投\quad [\href{https://arxiv.org/abs/2602.11877}{arXiv:2602.11877}]\quad \href{https://github.com/zhuchichi56/RouterXBench}{Code}}

\section{荣誉与奖项}
\begin{itemize}[itemsep=1pt, topsep=1pt]
\item \textbf{优秀毕业生}，南方科技大学致仁书院（2025.07）
\item \textbf{优秀毕业论文}，Top 5\%，南方科技大学（2025.07）
\item 优秀学生奖学金，南方科技大学（2024.10）
\end{itemize}
\section{技能}
\textbf{技术：}Python、PyTorch、Transformers、vLLM。\\
\textbf{英语：}IELTS 6.5。
\end{document}
